%%%%%%%%%%%%%%%%%%%%%%%%%%%%%%%%%%%%%%%%%%%%%%%%%%%%%
%			HLAVIČKA								%
%%%%%%%%%%%%%%%%%%%%%%%%%%%%%%%%%%%%%%%%%%%%%%%%%%%%%
\documentclass[openany,12pt]{memoir}

\usepackage[utf8]{inputenc} 
\usepackage[czech]{babel}
\usepackage[T1]{fontenc}
\usepackage[top=1.5cm, bottom=2cm, left=2cm, right=2cm]{geometry}  % --> NASTAVENÍ OKRAJŮ
\usepackage{fancyhdr}
\usepackage{graphicx}
%\usepackage{lmodern}
\usepackage{tgschola}  % --- Silnější font
\usepackage{xwatermark}
\usepackage{xcolor}
\usepackage{changepage}
\usepackage{newunicodechar}



%%%%%% Package na zpěvník
\usepackage[full]{leadsheets}%http://mirrors.nic.cz/tex-archive/macros/latex/contrib/leadsheets/leadsheets_en.pdf   --> dokumentace	
\definesongtitletemplate{empty}{} 
\setchords{
format = \bfseries,   %tučné akordy
minor = {mi},% 
input-notation = {german},%
output-notation = {german}%
}
\definesongtitletemplate{empty}{} 

\newlength{\drop}
\newwatermark[allpages,color=red!50,angle=0,scale=2, xpos=0,ypos=0]{\includegraphics[width=5cm]{obr/pozadi2.jpg}} %--> dvojka na pozadí


%%%% Vlastní příkazy
\newcounter{Slokočet}   %Automatické číslování slok
\newcommand{\mezera}{\vspace*{0.5cm}}   %Horizontální odsazení slok
\newcommand{\stred}{5.2cm}   %%% Na zarovnání slok doprostřed, pozn. automatičtější zarovnávání na střed nejde
\newcommand{\refren}{\mezera \noindent \textbf{R:} } %refrén
\newcommand{\sloka}{\mezera \noindent \addtocounter{Slokočet}{1} \arabic{Slokočet}. } 	%sloka, která se automaticky čísluje
\newcommand{\ssloka}{\mezera \noindent}  % vlastní číslo sloky
\newcommand{\carka}{,\:}
\newcommand{\m}[1]{\color{white}{#1}}  %Pro akordy
\newcommand{\ap}{'} %Pro apostrof


%%%%%%%%%%%%%%%%%%%%%%%%%%%%%%%%%%%%%%%%%%%%%%%%%%%%%%
%			NÁVOD									 %
%%%%%%%%%%%%%%%%%%%%%%%%%%%%%%%%%%%%%%%%%%%%%%%%%%%%%%
%1. Věci v hlavičce IGNOROVAT
%2. Píseň psát do prostoru mezi \begin{song} a \end{song}
%3. další řádek se značí dvěma odsazeníma (= dvakrát stisknout enter)
%4. \refren vždy na začátku refrenu a \sloka na začátek sloky (automaticky se čísluje)
%  \ = alt gr + q ; [/] = alt gr f/g ; {/} = alt gr + b/n; ^ = alt gr + 3 + mezera
%Cokoliv napíšete do ^{  } se bude brát jako akord
%když se toto bude dotýkat nějakého slova (nebude mezi tím a slovem mezera)
%tak se akordy zjeví nad slovem, ale pište to před slovo
%Když se to nedotýká slova, tak akord lítá ve vzduchu a vytiskne se větší mezera
%První možnost je asi preferovanější
%5. Akordy stačí psát jen do první sloky, když se nezmění -- kytaristi to zvládnou



%\usepackage{subfiles}
%</++>

\begin{document}

\pagestyle{empty}
%\addtocontents{toc}{\protect\thispagestyle{empty}}
%\tableofcontents \thispagestyle{empty}\newpage
\newgeometry{top=1.5cm, bottom = 0cm, left = 2cm, right = 2cm}

\newpage

%Pro kompilaci je důležitý, aby ten input a \end{document} zůstali na posledních
%dvou řádcích!!!!
\begin{song}{title=\centering De N\^{i}mes \\\normalsize Radůza  \vspace*{-0.3cm}}  %% sem se napíše jméno songu a autor
\moveright 1.7cm \vbox{      %Varianta č. 1  ---> Jeden sloupec zarovnaný na střed	

\begin{minipage}[t]{0.48\textwidth}\setlength{\parindent}{0.45cm}  %Varianta č. 2 --> Dva sloupce

\textbf{Ami\,\,Dmi\,\,G\,\,Emi\,\,Ami}
       
\sloka 
	^{Ami{\color{white}\_}}Očima uhla ^{Dmi}jsem

	^{G}a nebylo to ^{{\color{white}\_\_}Ami}studem,

	peřina pohla ^{Dmi}se

	^{G}nějak už ^{{\color{white}\_}Ami}bude.

	\phantom{.}

	Lino je studený

	ponožky a pak triko
   	
	a kabát přes denim
	
	a nic a ticho

\ssloka{\textbf{R_{1}}:}
	Cette chanson n’a ^{Ami}jamais été chantée

	comme il ^{Dmi}faut, c’est ça.

	Mais si tu ^{G}veux, nous tous les ^{Emi}deux
  
	pouvons le ^{Ami}faire, comme ça.

	\phantom{.}

	La tranquillité, c’est ce qu’on cherche
	
	on dit: ,,Puisse-t-elle venir!''
	
	On dit que tout est dificile

	mais, qu’ est-ce que ça veut dire.

\ssloka{\textbf{R_{2}}:}
	Zouvám si pár těžkejch bot
	
	si ve vesmíru můj pevnej bod
	
	a sukně kasám, je blízko brod
	
	a pryč je, co kdo zbod.

\end{minipage}\begin{minipage}[t]{0.48\textwidth}\setlength{\parindent}{0.45cm}\vspace*{0.55cm}  % V případě varianty č.2 jde odsud text do pravé části

\sloka
	Jen hlesnu od dveří
	
	víš, nelhala jsem dosud

	a jestli nevěříš,
	
	zkus špičku nosu.

	Okna jsou zamlžený,

	nádobí škopek plnej,

	mám hrdlo stažený
	
	hajej a nynej.

\ssloka{\textbf{R_{1}}:}

\ssloka{\textbf{R_{2}}:}

\phantom{.}

2x 
\textbf{Ami\,\,Dmi\,\,G\,\,Emi\,\,Ami}
       
\ssloka{\textbf{R_{2}}:}

\sloka
	Očima uhla jsem
	
	a nebylo to studem
	
	poznal si po hlase,

	jak mi teď bude.
	
	Lino je studený

	pod kafem stůl se kejvá
	
	un souvenir de N\^{i}mes
	
	tak už to bejvá.

\ssloka{\textbf{R_{1}}:}

\ssloka{\textbf{R_{2}}:}



\end{minipage}
}
\vspace*{2cm}
\moveright 1.7cm \vbox{      %Varianta č. 1  ---> Jeden sloupec zarovnaný na střed	

% Výslovnostní poznámka:\\
\phantom{.}

\begin{minipage}[t]{0.48\textwidth}\setlength{\parindent}{0.45cm}  %Varianta č. 2 --> Dva sloupce
\textbf{R_1}: 

  [Set šanson nežame ete šanté 

  kom il fo se sa 

  me si ty v ni tu led

  puvon le fér, kom sa.]

\end{minipage}\begin{minipage}[t]{0.48\textwidth}\setlength{\parindent}{0.45cm}%\vspace*{0.55cm}  % V případě varianty č.2 jde odsud text do pravé části
\phantom{.}

[Le trankilité se s kon šerš

  on dy pui tel venir

  on dy k tu e dyfisil

  me kes ke sa v dire.]
\end{minipage}\\\\\\
% Texty překládat nebudeme. Věříme, že o významu jakékoliv písničky by člověk\\ měl přemýšlet především sám.
}
\setcounter{Slokočet}{0}
\end{song}

\end{document}
